% from ega3-1-fr.tex
\oldpage[III]{103}
soit bijectif. En d'autres termes, l'homomorphisme canonique
\[
\alpha:\mathcal M\longrightarrow\bbGamma_*(\shProj(\mathcal M))
\]
est un \textbf{(TN)}-isomorphisme.
\end{theorem}

\begin{proof}
On peut se borner au cas où $\mathcal M$ est un $\mathcal S$-module de
type fini (II, 3.4.2). Comme $Y$ est quasi-compact, il existe un entier
$d>0$ tel que $\mathcal S^{(d)}$ soit engendrée par le
$\mathcal O_Y$-Module quasi-cohérent $\mathcal S_d$, ce dernier étant de
type fini (II, 3.1.10), donc cohérent puisque $Y$ est noethérien.
Remarquons maintenant que $\mathcal M$ est somme directe des
$\mathcal M^{(d,k)}$ pour $0\leq k<d$ et que chacun des
$\mathcal M^{(d,k)}$ est un $\mathcal S^{(d)}$-Module quasi-cohérent de
type fini, ainsi qu'il résulte de (II, 2.1.6, (iii)), la question étant
locale sur $Y$. Or, il suffit évidemment de prouver que chacun des
homomorphismes canoniques
\[
\alpha:\mathcal M^{(d,k)}
\longrightarrow\bbGamma_*((\shProj(\mathcal M))^{(d,k)})
\]
est un \textbf{(TN)}-isomorphisme. Compte tenu de (II, 8.14.13), et
notamment du diagramme (8.14.13.4), on voit qu'on est ramené à prouver le
théorème lorsque $\mathcal S$ est engendrée par $\mathcal S_1$ et que
$\mathcal S_1$ est un $\mathcal O_Y$-Module cohérent.

Comme $Y$ est noethérien, le même raisonnement qu'au début de
(\hyperref[III.2.2.2-fr]{2.2.2}) montre qu'on peut se borner au cas où
$Y=\operatorname{Spec}(A)$, $\mathcal S=\widetilde S$,
$\mathcal M=\widetilde M$, $A$ étant un anneau noethérien, $S_1$ un
$A$-module de type fini et $M$ un $S$-module gradué de type fini.
Montrons qu'il suffit alors de prouver le théorème lorsque $M=S$. En effet,
dans le cas général, on a une suite exacte
\[
L'\longrightarrow L\longrightarrow M\longrightarrow0,
\]
où $L$ et $L'$ sont des sommes directes de modules gradués de la forme
$S(m)$. Si le résultat est vrai pour $M=S$, il l'est aussi pour $M=S(m)$,
donc pour $L$ et $L'$. Considérons alors le diagramme commutatif
\[
\xymatrix{
\widetilde {L'_n}\ar[r]\ar[d]_{\alpha_n} &
\widetilde {L_n}\ar[r]\ar[d]^{\alpha_n} &
\widetilde {M_n}\ar[r]\ar[d]^{\alpha_n} & 0\\
q_*(\widetilde {L'}(n))\ar[r] &
q_*(\widetilde L(n))\ar[r] &
q_*(\widetilde M(n))\ar[r] & 0 .
}
\]
La deuxième ligne est exacte en vertu de
(\hyperref[III.2.2.3-fr]{2.2.3}), dès que $n$ est assez grand~; comme il
en est de même de la première et que les deux flèches verticales de gauche
sont des isomorphismes, il en est de même de la troisième.

Cela étant, pour prouver le théorème lorsque $M=S$, supposons d'abord que
$S=A[T_0,\ldots,T_r]$ ($T_i$ indéterminées)~; dans ce cas, notre assertion
n'est autre que (\hyperref[III.2.1.11-fr]{2.1.11}, (ii)). Dans le cas
général, $S$ s'identifie à un quotient d'un anneau
$S'=A[T_0,\ldots,T_r]$ par un idéal gradué, donc $X$ à un sous-schéma
fermé de $X'=\mathbf P_A^r$ (II, 2.9.2). Si $j$ est l'injection canonique
$X\to X'$, $j_*(\widetilde S(n))$ n'est autre que le
$\mathcal O_{X'}$-Module $(\shProj(\widetilde S))(n)$ où $S$ est considéré
comme un $S'$-module gradué~; cela résulte en effet aussitôt de
(II, 2.8.7). Comme $j_*(\widetilde S(n))$ est un
$\mathcal O_{X'}$-Module vérifiant \textbf{(TF)}, l'homomorphisme canonique
$\alpha_n:S_n\to\Gamma(X',j_*(\widetilde S(n)))$ est bijectif pour $n$
assez grand, en vertu de ce qui précède~; cela achève la démonstration,
puisque
$\Gamma(X',j_*(\widetilde S(n)))=\Gamma(X,\widetilde S(n))$.
\end{proof}

